%%%%%%%%%%%%%%%%%%%%%%%%%%%%%%%%%%%%%%%%%
% Specjalna strona pracy ze streszczeniem i abstractem w j. angielskim
% Szablon pracy dyplomowej
% Wydział Informatyki 
% Zachodniopomorski Uniwersytet Technologiczny w Szczecinie
% autor Joanna Kołodziejczyk (jkolodziejczyk@zut.edu.pl)
% Bardzo wczesnym pierwowzorem szablonu był
% The Legrand Orange Book
% Version 5.0 (29/05/2025)
%
% Modifications to LOB assigned by %JK
%%%%%%%%%%%%%%%%%%%%%%%%%%%%%%%%%%%%%%%%%


\begin{center}
\noindent {{\color{blueZUT}\Large\sffamily  {Streszczenie}}}\\[1cm] 
\end{center}

   Przedmiotem niniejszej pracy jest projekt i implementacja systemu informatycznego ,,Household Manager'' wspomagającego zarządzanie gospodarstwem domowym oraz zintegrowanymi urządzeniami inteligentnego domu (smart home). Aplikacja rozwiązuje problemy związane z podziałem obowiązków, rozliczaniem wspólnych wydatków oraz koordynacją działań między domownikami.

System składa się z aplikacji mobilnej stworzonej w technologii React Native oraz serwera backendowego napisanego w języku Go, wykorzystującego bazę danych PostgreSQL. Kluczowe funkcjonalności obejmują: zarządzanie wirtualnymi domami, moduł sterowania urządzeniami smart home, system przydzielania zadań, moduł finansowy z automatycznym odczytywaniem danych z rachunków (OCR), wspólne listy zakupów oraz ankiety decyzyjne. Zastosowanie nowoczesnych technologii pozwoliło na stworzenie wydajnego, skalowalnego i intuicyjnego w obsłudze rozwiązania, dostępnego na platformy Android i iOS.

\vspace{10pt}
\noindent{\bf słowa kluczowe:} zarządzanie gospodarstwem domowym, inteligentny dom, aplikacja mobilna, podział obowiązków, wspólne wydatki, React Native, Go

\vfill

\begin{center}
\noindent {{\color{blueZUT}\Large\sffamily {Abstract}}}\\[1cm] 
\end{center}

The subject of this thesis is the design and implementation of the “Household Manager” IT system, which supports household management and integrated smart home devices. The application solves problems related to the division of responsibilities, settlement of joint expenses, and coordination of activities between household members.

The system consists of a mobile application created in React Native technology and a backend server written in Go, using a PostgreSQL database. Key features include: virtual home management, a smart home device control module, a task assignment system, a financial module with automatic bill reading (OCR), shared shopping lists, and decision-making surveys. The use of modern technologies has allowed us to create an efficient, scalable, and intuitive solution, available on Android and iOS platforms.

\vspace{10pt}
\noindent{\bf keywords:} household management, smart home, mobile application, task management, shared expenses, React Native, Go
